\section{Literature}\label{sec:literature}
Our work relates to three main strands of the literature.

\subsection{School Choice and Student Outcomes}
A large empirical literature examines how expanding school choice---i.e., relaxing traditional neighborhood-based assignment policies and allowing families to select among multiple schools, including magnet and voucher programs, open-enrollment initiatives, and inter-district choice---affects student outcomes. Early work focused on cross-district or cross-municipality variation in competitive pressure, with mixed findings on achievement, school quality and productivity~\citep{hoxby2000, hoxby2003, hsieh2006, rothstein2007}. Other studies exploit lotteries or oversubscription rules to estimate the individual effects of attending preferred schools~\citep{cullen2006, abdulkadiroglu2011, deming2014, muralidharan2015}. These papers typically evaluate short-run impacts on test scores, peer composition, or satisfaction, and do not study how centralized assignment shapes longer-term educational trajectories.
The study most closely related to ours is \citet{camposkearns2024}, who analyze Los Angeles’s Zones of Choice (ZOC), a district-level reform that created small local choice markets while retaining neighborhood boundaries elsewhere. The authors show that student outcomes in ZOC markets increased markedly, narrowing achievement and college enrollment gaps between ZOC neighborhoods and the rest of the district. Furthermore, the authors identify competition-induced improvements in school quality as the central mechanism behind the effects of the reform. 




\subsection{Centralized School Choice.}
A parallel literature studies the design and consequences of \emph{centralized} school choice systems, in which families submit rank-ordered preference lists and a central clearinghouse assigns students to schools using a matching algorithm. Starting with \citet{abdulkadiroglu2003}, who formalize school choice as a mechanism design problem and compare the properties of the two most widely used assignment rules—\emph{deferred acceptance} and \emph{immediate acceptance}—this literature has developed along several complementary lines. 
First, a large theoretical literature studies the incentive, efficiency, and stability properties of alternative matching mechanisms and the strategic behavior they induce \citep{abdulkadiroglu2003,kesten2010,pathak2017}. Second, a set of papers evaluates the implementation of centralized assignment systems in large urban school districts, documenting how changes in the mechanism affect equilibrium behavior, strategic reporting of preferences, and assignment outcomes \citep{abdulkadiroglu2005,abdulkadiroglu2009}. Third, a growing empirical literature uses administrative data from centralized clearinghouses to estimate students’ preferences over schools and to evaluate welfare and distributional consequences of school choice systems \citep{hastings2008,agarwal2018}. Recent work also studies how centralized assignment reforms affect broader outcomes such as access to high-quality schools, school segregation, and students’ educational trajectories \citep{neilson2021, gazmuri2024}. 

Within this literature, the paper most closely related to ours is \cite{urzua2023}, which studies Chile’s transition to a centralized school‑choice system and employs a similar staggered difference‑in‑differences strategy to estimate the reform's effects on schools' socioeconomic composition and within‑school segregation. Our analysis complements theirs by tracing the reform's downstream effects on college‑admissions outcomes, asking whether changes in assignment and school composition altered students' application behavior and admission outcomes.

% Urzúa et al. find no overall improvement in low‑SES representation and document increased within‑school segregation in areas with high residential segregation or many private schools, driven in part by high‑SES families’ migration to private institutions. This addition highlights why your paper’s focus on college admissions is a natural and important complement to their work.

\subsection{Centralized College Admissions.}
Finally, our work is also related to a literature studying the drivers of outcomes in \emph{centralized college admissions systems}. 
One strand of this literature focuses on the \emph{design of the assignment mechanism}. In particular, this work studies how different elements of the allocation process---including the algorithm used~\citep{Calsamiglia2020}, exam-retaking and re-application policies~\citep{larroucau2025, rios2026}, tie-breaking rules~\citep{abdulkadiroglu2009, Arnosti2015, Ashlagi2019}, affirmative action policies~\citep{rios2021improving,Baswana2019,otero23}, the timing of allocations~\citep{Luflade2017, GongLiang2022}, and the aftermarket~\citep{Kapor2024}---affect the theoretical properties of the mechanism, students' strategic behavior, and equilibrium allocation outcomes~\citep{chenkoesten2013,fu2014,neilson2021}. 

Another strand studies the role of information frictions and beliefs in shaping students’ application behavior. In many centralized systems, students lack complete information about program characteristics, hold biased beliefs about their admission probabilities, or are unaware of alternative programs that may be a better fit. As a result, they may submit suboptimal preference lists, mis‑rank programs, or fail to update their plans in response to new information. Empirical work has documented how targeted information and advising can improve match quality: \citet{hastings2008} and \citet{agarwal2018} show that information interventions improve application behavior and outcomes; \citet{Narita2019} exploits a staggered rollout of school choice information to estimate learning dynamics and inertia; and \citet{larroucau2024} studies how beliefs about admission probabilities shape strategic ranking behavior across heterogeneous students. More generally, this literature highlights how informational constraints interact with the mechanism to shape final match quality and downstream outcomes.

Within this strand of the literature, the most closely related to ours are papers that use quasi‑experimental variation in assignment rules to document effects on downstream educational outcomes. \citet{terrier2024} exploit England's 2008 nationwide ban on the Immediate Acceptance (IA) mechanism, which forced local authorities to switch to Deferred Acceptance (DA), and use a difference‑in‑differences design to show that the transition to DA reduced low‑SES students' access to high value‑added and selective schools and widened socioeconomic gaps in school quality. They attribute this to a competition‑for‑top‑schools effect: under IA, high‑SES families had strategically avoided oversubscribed schools, but once DA removed those strategic incentives they flooded back, crowding out disadvantaged applicants — an effect concentrated in competitive local authorities with many selective schools. \citet{kaporneilsonzimmerman2020} similarly find that when parents hold incorrect beliefs about admission chances, DA can dominate IA in welfare terms, with gains concentrated among low‑SES families, suggesting that the information environment is a key moderator of mechanism performance. Moving to longer‑run outcomes, \citet{cullenjacobleviit2006} exploit lottery variation in Chicago's school choice program and find little effect of winning a choice lottery on academic achievement, while \citet{demingetal2014} use a similar design in Charlotte‑Mecklenburg and find positive effects on postsecondary attainment for low‑income students when assignment mechanisms successfully place them in higher‑quality schools. Taken together, these papers suggest that centralizing school choice reshapes the distribution of school quality across socioeconomic groups and can affect students' long‑run trajectories, but that the direction and magnitude of effects depend critically on the degree of competition for desirable seats and on whether selective admissions criteria interact with strategic behavior in ways that systematically disadvantage less sophisticated families.







% Together, this body of work highlights how the design of centralized matching mechanisms shapes both the strategic environment faced by families and the allocation of students across schools, with important implications for equity and efficiency in education markets.


\subsection{Contribution}
Our work contributes to and extends these literatures in two important ways. First, much of the existing empirical work on school choice evaluates policy changes implemented at the district or local level that expand families' ability to choose among schools. In contrast, the reform we study represents a nationwide institutional transformation: the introduction of the \emph{Sistema de Admisión Escolar} (SAE), which replaced decentralized, school-level admissions processes with a unified centralized clearinghouse based on a deferred-acceptance algorithm. This reform eliminated discretionary selection practices by schools and standardized admissions rules across all publicly funded institutions. As a result, it provides a unique setting to study how a large-scale transition to centralized assignment affects student behavior and outcomes.

Second, our paper examines a different set of outcomes than most of the existing literature. While prior work on school choice and centralized assignment primarily focuses on short-run outcomes such as test scores, peer composition, or access to preferred schools, our data allow us to link administrative records across secondary and postsecondary education. This linkage enables us to study how exposure to centralized school assignment during secondary education affects students’ subsequent participation and outcomes in a separate nationwide centralized college admissions system.

More broadly, the literature on centralized college admissions emphasizes that outcomes depend critically on institutional design and students’ application behavior within the system. In contrast, our analysis highlights how experiences \emph{prior} to college admissions can shape behavior in centralized matching environments. Specifically, we study whether participation in a centralized school choice clearinghouse during secondary education affects students’ application, admission, and enrollment outcomes when they later participate in centralized college admissions. By linking these two centralized markets, our paper shows how earlier exposure to centralized choice institutions can influence behavior and outcomes in subsequent matching environments.

% Finally, the mechanisms differ. In ZOC, competitive pressures among schools are the primary driver of improvements: schools respond to expanded choice by increasing quality. In contrast, SAE’s effects arise from eliminating selective screening, reducing information frictions, and standardizing assignment rules. 

